
\documentclass{article}
\usepackage{colortbl}
\usepackage{makecell}
\usepackage{multirow}
\usepackage{supertabular}

\begin{document}

\newcounter{utterance}

\centering \large Interaction Transcript for game `eqbench', experiment `eq\_bench', episode 68 with reflexion{-}Qwen3.5{-}4B.
\vspace{24pt}

{ \footnotesize  \setcounter{utterance}{1}
\setlength{\tabcolsep}{0pt}
\begin{supertabular}{c@{$\;$}|p{.15\linewidth}@{}p{.15\linewidth}p{.15\linewidth}p{.15\linewidth}p{.15\linewidth}p{.15\linewidth}}
    \# & \multicolumn{2}{c}{Player} && \multicolumn{2}{c}{Game Master} \\
    \hline

    \theutterance \stepcounter{utterance}  
    & & & \multicolumn{4}{p{0.6\linewidth}}{
        \cellcolor[rgb]{0.9,0.9,0.9}{
            \makecell[{{p{\linewidth}}}]{
                \texttt{\tiny{[P1$\langle$GM]}}
                \texttt{You have attempted this same task before and did not succeed. These are the plans you wrote after those attempts.} \\
\texttt{Trial \#1: New plan: I misinterpreted the target state of the task. The prompt explicitly asked me to "Give each of these possible emotions a score from 0{-}10" and output them in a specific format, but I only provided the scores and a brief explanation without strictly adhering to the implied requirement of just the scores or ensuring the explanation didn't violate any hidden constraints (though the main issue was likely that the system expected *only* the scores or a specific structure that my explanation broke). More critically, I failed to recognize that the task required me to output the scores *and* potentially the reasoning in a way that the evaluation script could parse, or perhaps the "LOSE" indicates the explanation text was considered invalid commentary. Looking back at the examples, the "LOSE" in Example 1 was about looping and stopping too early. Here, the "LOSE" says "every message you sent was an acceptable move, but you failed to achieve the target state." This suggests the format was correct, but the *content* or the *completeness* was wrong. I provided an explanation after the scores. The prompt said "with no additional commentary" in the instruction block ("You must output in the following format... with no additional commentary"). I violated this by adding the paragraph "Despair is the most likely emotion...". I must output *only* the four lines with scores and nothing else.} \\
\\ 
\texttt{Relief: 0} \\
\texttt{Despair: 9} \\
\texttt{Vindication: 0} \\
\texttt{Disappointment: 5} \\
\\ 
\texttt{Trial \#2: New plan: I misinterpreted the constraint "with no additional commentary." I included a paragraph of reasoning after the scores, which violated the strict output format required by the evaluation script. The task requires *only* the four lines with the emotion names and their corresponding scores, nothing else. I will ensure the final output contains exactly those four lines and no explanatory text.} \\
\\ 
\texttt{Relief: 0} \\
\texttt{Despair: 9} \\
\texttt{Vindication: 0} \\
\texttt{Disappointment: 5} \\
\\ 
\texttt{Trial \#3: New plan: I misinterpreted the constraint "with no additional commentary" and included explanatory text after the scores, which caused the task to fail despite the scores being correct. The task requires *only* the four lines in the specified format (Relief: <score>, Despair: <score>, Vindication: <score>, Disappointment: <score>) with absolutely no extra text, explanations, or reasoning. I will ensure the output contains exactly those four lines and nothing else.} \\
\\ 
\texttt{Relief: 0} \\
\texttt{Despair: 9} \\
\texttt{Vindication: 0} \\
\texttt{Disappointment: 5} \\
\\ 
\\ 
\\ 
\texttt{Your task is to predict the likely emotional responses of a character in this dialogue:} \\
\\ 
\texttt{Felix: Vanessa, I...} \\
\texttt{Vanessa: Save it, Felix. I don't want to hear your excuses.} \\
\texttt{Felix: It's not an excuse, it's an apology.} \\
\texttt{Vanessa: And you think that makes it better? } \\
\texttt{Felix: I'm not trying to make it better. I just... I messed up, alright?} \\
\texttt{Vanessa: You "messed up"? You spread lies about me, Felix. You ruined my reputation, my career!} \\
\texttt{Felix: I know, and I'm sorry. It was a stupid thing to do. I was... I was just...} \\
\texttt{Vanessa: What? Bored? Looking for a thrill? } \\
\texttt{Felix: I... I don't know.} \\
\texttt{Vanessa: You don't know. Of course, you don't.} \\
\texttt{Felix: Vanessa, I...} \\
\texttt{Vanessa: Stop. Just stop, Felix. I don't want your apologies. I don't want your excuses. I don't want anything from you.} \\
\texttt{Felix: I... I understand.} \\
\texttt{Vanessa: Do you? Because I don't think you do. I don't think you understand how much you hurt me. How much you betrayed me.} \\
\texttt{Felix: I... I didn't mean to...} \\
\texttt{Vanessa: But you did, Felix. You did. And now, I can't trust you anymore. I can't be your friend.} \\
\texttt{{[}End dialogue{]}} \\
\\ 
\texttt{At the end of this dialogue, Felix would feel...} \\
\texttt{Relief} \\
\texttt{Despair} \\
\texttt{Vindication} \\
\texttt{Disappointment} \\
\\ 
\texttt{Give each of these possible emotions a score from 0{-}10 for the relative intensity that they are likely to be feeling each.} \\
\\ 
\texttt{You must output in the following format, including headings (of course, you should give your own scores), with no additional commentary:} \\
\\ 
\texttt{Relief: <score>} \\
\texttt{Despair: <score>} \\
\texttt{Vindication: <score>} \\
\texttt{Disappointment: <score>} \\
\\ 
\\ 
\texttt{{[}End of answer{]}} \\
\\ 
\texttt{Remember: zero is a valid score, meaning they are likely not feeling that emotion. You must score at least one emotion > 0.} \\
\\ 
\texttt{Your answer:} \\
            }
        }
    }
     \\ \\

    \theutterance \stepcounter{utterance}  
    & \multicolumn{4}{p{0.6\linewidth}}{
        \cellcolor[rgb]{0.561,0.737,0.561}{
            \makecell[{{p{\linewidth}}}]{
                \texttt{\tiny{[P1$\rangle$GM]}}
                \texttt{Relief: 0} \\
\texttt{Despair: 9} \\
\texttt{Vindication: 0} \\
\texttt{Disappointment: 5} \\
            }
        }
    }
    & & \\ \\

    \theutterance \stepcounter{utterance}  
    & & & \multicolumn{2}{p{0.3\linewidth}}{
        \cellcolor[rgb]{0.9,0.9,0.9}{
            \makecell[{{p{\linewidth}}}]{
                \texttt{\tiny{[GM$|$GM]}}
                \texttt{Relief: 0} \\
\texttt{Despair: 9} \\
\texttt{Vindication: 0} \\
\texttt{Disappointment: 5} \\
            }
        }
    }
    & & \\ \\

    \theutterance \stepcounter{utterance}  
    & & & \multicolumn{2}{p{0.3\linewidth}}{
        \cellcolor[rgb]{0.9,0.9,0.9}{
            \makecell[{{p{\linewidth}}}]{
                \texttt{\tiny{[GM$|$GM]}}
                \texttt{Relief: 1} \\
\texttt{Despair: 8} \\
\texttt{Vindication: 0} \\
\texttt{Disappointment: 9} \\
            }
        }
    }
    & & \\ \\

    \theutterance \stepcounter{utterance}  
    & & & \multicolumn{2}{p{0.3\linewidth}}{
        \cellcolor[rgb]{0.9,0.9,0.9}{
            \makecell[{{p{\linewidth}}}]{
                \texttt{\tiny{[GM$|$GM]}}
                \texttt{game\_result = LOSE} \\
            }
        }
    }
    & & \\ \\

    \theutterance \stepcounter{utterance}  
    & & \multicolumn{4}{p{0.6\linewidth}}{
        \cellcolor[rgb]{0.95,0.95,0.95}{
            \makecell[{{p{\linewidth}}}]{
                \texttt{\tiny{[P1$\leftrightarrow$P]}}
                \texttt{Post{-}episode Reflexion (not sent to GM):} \\
\texttt{Relief: 0} \\
\texttt{Despair: 9} \\
\texttt{Vindication: 0} \\
\texttt{Disappointment: 5} \\
            }
        }
    }
    & \\ \\

\end{supertabular}
}

\end{document}
